\begin{thebibliography}{99}
\bibitem{addario2014growing} L. Addario-Berry - “Growing random 3-connected maps” , Electron. Comm. Probab. 19 (2014), article ID 54, 12 pages
\bibitem{MR3035742} R. Abraham, J.-F. Delmas \& P. Hoscheit - “A note on the Gromov-Hausdorff-Prokhorov distance between (locally) compact metric measure spaces” , Electron. J. Probab. 18 (2013), article ID 14, 21 pages
\bibitem{Ald91a} D. Aldous - “The continuum random tree. I” , Ann. Probab. 19 (1991) no. 1, p. 1-28
\bibitem{Ald91} D. Aldous - “The continuum random tree. II. An overview” , in Stochastic analysis (Durham, 1990) , London Math. Soc. Lecture Note Ser. , vol. 167 , Cambridge Univ. Press, Cambridge, 1991, p. 23-70
\bibitem{Ald93} D. Aldous - “The continuum random tree. III” , Ann. Probab. 21 (1993) no. 1, p. 248-289
\bibitem{AthreayLorWinter} S. Athreya, W. Löhr \& A. Winter - “The gap between Gromov-vague and Gromov–Hausdorff-vague topology” , Stochastic Process. Appl. 126 (2016) no. 9, p. 2527-2553
\bibitem{bacher2017efficient} A. Bacher, O. Bodini \& A. Jacquot - “Efficient random sampling of binary and unary-binary trees via holonomic equations” , Theoret. Comput. Sci. 695 (2017), p. 42-53
\bibitem{bertoin2002self} J. Bertoin - “Self-similar fragmentations” , Ann. Inst. H. Poincaré Probab. Statist. 38 (2002) no. 3, p. 319-340
\bibitem{billingsley1965} P. Billingsley - Ergodic theory and information , Wiley Series in Probab. and Math. Stat. , Wiley, 1965
\bibitem{BK14} J. Bertoin \& I. Kortchemski - “Self-similar scaling limits of Markov chains on the positive integers” , Ann. Appl. Probab. 26 (2016) no. 4, p. 2556-2595
\bibitem{CFT25} N. Curien, W. Fleurat \& A. Twigt - “Growing self-similar Markov trees” , 2025, work in progress
\bibitem{CurienHaas} N. Curien \& B. Haas - “The stable trees are nested” , Probab. Theory Relat. Fields 157 (2013) no. 3-4, p. 847-883
\bibitem{CLGharmonic} N. Curien \& J.-F. Le Gall - “The harmonic measure of balls in random trees” , Ann. Probab. 45 (2017) no. 1, p. 147-209
\bibitem{Carmona/Petit/Yor} P. Carmona, F. Petit \& M. Yor - “On the distribution and asymptotic results for exponential functionals of Lévy processes” , in Exponential functionals and principal values related to Brownian motion , Bibl. Rev. Mat. Iberoamericana , Revista Matemática Iberoamericana, Madrid, 1997, p. 73-130
\bibitem{CS20} A. Caraceni \& A. Stauffer - “Polynomial mixing time of edge flips on quadrangulations” , Probab. Theory Relat. Fields 176 (2020) no. 1-2, p. 35-76
\bibitem{CS23} A. Caraceni \& A. Stauffer - “Growing uniform planar maps face by face” , Random Structures Algorithms 63 (2023) no. 4, p. 942-967
\bibitem{duquesne2005probabilistic} T. Duquesne \& J.-F. Le Gall - “Probabilistic and fractal aspects of Lévy trees” , Probab. Theory Relat. Fields 131 (2005) no. 4, p. 553-603
\bibitem{duquesne2006hausdorff} T. Duquesne \& J.-F. Le Gall - “The Hausdorff measure of stable trees” , Alea 1 (2006), p. 393-415
\bibitem{MR2221786} S. N. Evans, J. Pitman \& A. Winter - “Rayleigh processes, real trees, and root growth with re-grafting” , Probab. Theory Relat. Fields 134 (2006) no. 1, p. 81-126
\bibitem{EW06} S. N. Evans \& A. Winter - “Subtree prune and regraft: a reversible tree-valued Markov process” , Ann. Probab. 34 (2006) no. 3, p. 81-126
\bibitem{FPRW18} N. Forman, S. Pal, D. Rizzolo \& M. Winkel - “Aldous diffusion I: a projective system of continuum $k$ -tree evolutions” , 2018
\bibitem{FPRW20} N. Forman, S. Pal, D. Rizzolo \& M. Winkel - “Projections of the Aldous chain on binary trees: intertwining and consistency” , Random Structures Algorithms 57 (2020) no. 3, p. 745-769
\bibitem{forman2023aldous} N. Forman, S. Pal, D. Rizzolo \& M. Winkel - “The Aldous diffusion: a stationary evolution of the Brownian CRT” , 2023
\bibitem{goldschmidt2019parking} C. Goldschmidt \& M. Przykucki - “Parking on a random tree” , Combin. Probab. Comput. 28 (2019) no. 1, p. 23-45
\bibitem{HM04} B. Haas \& G. Miermont - “The genealogy of self-similar fragmentations with negative index as a continuum random tree” , Electron. J. Probab. 9(4) (2004), p. 57-97
\bibitem{HM12} B. Haas \& G. Miermont - “Scaling limits of Markov branching trees with applications to Galton-Watson and random unordered trees” , Ann. Probab. 40 (2012) no. 6, p. 2589-2666
\bibitem{HS14} B. Haas \& R. Stephenson - “Scaling limits of $k$ -ary growing trees” , Ann. Inst. H. Poincaré Probab. Statist. 51 (2015) no. 4, p. 1314-1341
\bibitem{khezeli2020metrization} A. Khezeli - “Metrization of the Gromov–Hausdorff(-Prokhorov) topology for boundedly-compact metric spaces” , Stochastic Process. Appl. 130 (2020) no. 6, p. 3842-3864
\bibitem{khezeli2023unified} A. Khezeli - “A unified framework for generalizing the Gromov-Hausdorff metric” , Probab. Surv. 20 (2023), p. 837-896
\bibitem{Lamp62} J. Lamperti - “Semi-stable stochastic processes” , Trans. Amer. Math. Soc. 104 (1962), p. 62-78
\bibitem{lawler1993discrete} G. F. Lawler - “A discrete analogue of a theorem of Makarov” , Combin. Probab. Comput. 2 (1993) no. 2, p. 181-199
\bibitem{LG2005} J.-F. Le Gall - “Random trees and applications” , Probab. Surv. 2 (2005), p. 245-311
\bibitem{lohr2020aldous} W. Löhr, L. Mytnik \& A. Winter - “The Aldous chain on cladograms in the diffusion limit” , Ann. Probab. 48 (2020) no. 5, p. 2565-2590
\bibitem{lyons1995ergodic} R. Lyons, R. Pemantle \& Y. Peres - “Ergodic theory on Galton–Watson trees: speed of random walk and dimension of harmonic measure” , Ergodic Theory Dynam. Systems 15 (1995) no. 3, p. 593-619
\bibitem{lyons2008growth} R. Lyons, R. Peled \& O. Schramm - “Growth of the number of spanning trees of the Erdős–Rényi giant component” , Combin. Probab. Comput. 17 (2008) no. 5, p. 711-726
\bibitem{LW04} M. Luczak \& P. Winkler - “Building uniformly random subtrees” , Random Structures Algorithms 24 (2004) no. 4, p. 420-443
\bibitem{Mar03} P. Marchal - “Constructing a sequence of random walks strongly converging to Brownian motion” , in Discrete random walks (Paris, 2003) , Discrete Math. Theor. Comput. Sci. Proc., AC , Assoc. Discrete Math. Theor. Comput. Sci., Nancy, 2003, p. 181-190
\bibitem{marckert2023growing} J.-F. Marckert - “Growing random uniform d-ary trees” , Ann. Comb. 27 (2023) no. 1, p. 51-66
\bibitem{mcshine1997mixing} L. McShine \& P. Tetali - “On the mixing time of the triangulation walk and other Catalan structures” , in Randomization methods in algorithm design (Princeton, NJ, 1997) , DIMACS Ser. Discrete Math. Theoret. Comput. Sci. , vol. 43 , American Mathematical Society, Providence, RI, 1999, p. 147-160
\bibitem{PitmanWinkel2013} J. Pitman \& M. Winkel - “Regenerative tree growth: Markovian embedding of fragmenters, bifurcators, and bead splitting processes” , Ann. Probab. 43 (2015) no. 5, p. 2611-2646
\bibitem{Rem85} J.-L. Rémy - “Un procédé itératif de dénombrement d’arbres binaires et son application à leur génération aléatoire” , RAIRO Inform. Théor. 19 (1985) no. 2, p. 179-195
\bibitem{Steph13} R. Stephenson - “General fragmentation trees” , Electron. J. Probab. 18(101) (2013), p. 1-45
\end{thebibliography}
